% Colours of the presentation template, defined here as well so that the TikZ diagrams also compile in the cheat sheet.
\definecolor{nyupurple}{RGB}{90, 35, 140}
\definecolor{lightnyupurple}{RGB}{159, 120, 184}

% GRU diagram, built in stages: #1 = 1 draws the update gate, #2 = 1 draws the reset gate.
\providecommand{\grudiagram}[2]{%
\adjustbox{max width=\linewidth, max totalheight=0.42\textheight}{%
\begin{tikzpicture}[>=latex, thick, every node/.style={font=\small}, st/.style={draw, circle, minimum size=9mm, inner sep=1pt}, gt/.style={draw, rounded corners, fill=orange!25, minimum height=7mm, align=center, inner sep=3pt}]
\node[st] (hp) at (0,2) {$h_{t-1}$};
\node[st] (xt) at (0,-0.25) {$x_t$};
\node[draw, rounded corners, fill=lightnyupurple!30, minimum width=26mm, minimum height=12mm, align=center] (cd) at (5,0) {candidate $\tilde{h}_t$\\$\tanh$};
\node[draw, rounded corners, fill=white, minimum width=11mm, minimum height=7mm] (bl) at (9,2) {mix};
\node[st] (ht) at (11.5,2) {$h_t$};
\node[font=\scriptsize, gray, above] at (hp.north) {previous state};
\node[font=\scriptsize, gray, below] at (xt.south) {input};
\node[font=\scriptsize, gray, above] at (ht.north) {new state};
\draw[->, line width=1.6pt, lightnyupurple!70!black] (hp) -- node[above, font=\scriptsize, black] {old memory, passed on} (bl);
\fill (2,2) circle (2pt);
\draw[->] (xt.east) -- (3.7,-0.25);
\ifnum#2=1
\node[gt] (rg) at (2,1.05) {reset\\$r_t$};
\draw (2,2) -- (rg.north);
\draw[->] (rg.south) -- (2,0.25) -- (3.7,0.25);
\else
\draw[->] (2,2) -- (2,0.25) -- (3.7,0.25);
\fi
\draw[->] (cd.east) -| node[pos=0.25, below, font=\scriptsize] {new content} (bl.south);
\ifnum#1=1
\node[gt] (zg) at (9,3.5) {update $z_t$};
\draw[->] (zg.south) -- (bl.north);
\fi
\draw[->] (bl) -- (ht);
\end{tikzpicture}}}

%%%%%%%%%%%%%%%%%%%%%%%%%%%%%%%%%%%%%%%%%%%%%%%%%%%
\begin{frame}
  \begin{center}
    {\Large Gated Recurrent Unit (GRU)}
  \end{center}
\end{frame}


%%%%%%%%%%%%%%%%%%%%%%%%%%%%%%%%%%%%%%%%%%%%%%%%%%%
\begin{frame}[fragile] \frametitle{Why Gates?}
\begin{itemize}
\item A plain RNN rewrites its whole state at every step: $h_t = \tanh(W_{xh}x_t + W_{hh}h_{t-1} + b_h)$
\item Nothing is copied unchanged: every old value passes through weights and a $\tanh$, step after step
\item So it cannot hold a value until it is needed, or ignore a useless reading such as a sensor glitch
\item A \textbf{gate} is a number between 0 and 1 for every memory slot (one component of the state vector)
  \begin{itemize}
  \item It comes from a sigmoid, a curve that squeezes any number into the range 0 to 1
  \item 0 blocks, 1 lets everything through, and in between lets a share through
  \end{itemize}
\item The GRU (Cho et al., 2014) uses gates so the cell can choose, step by step, between keeping and rewriting its memory
\end{itemize}
\begin{block}{Intuition}
A gate works like a tap on a pipe. Fully closed, nothing passes; fully open, everything passes. The network learns how far to open each tap at every step.
\end{block}
\end{frame}


%%%%%%%%%%%%%%%%%%%%%%%%%%%%%%%%%%%%%%%%%%%%%%%%%%%
\begin{frame}[fragile] \frametitle{The GRU Main Path}
\begin{center}
\grudiagram{0}{0}
\end{center}
\begin{itemize}
\item Same two inputs as a plain RNN: the old state $h_{t-1}$ and the new input $x_t$
\item The cell proposes a \textbf{candidate} $\tilde{h}_t$: what the new state could be
\item The new state is a \textbf{mix} of the old memory and the candidate
\item The old memory reaches the mix directly, with no weight matrix and no $\tanh$ on the way: a short path
\end{itemize}
\end{frame}


%%%%%%%%%%%%%%%%%%%%%%%%%%%%%%%%%%%%%%%%%%%%%%%%%%%
\begin{frame}[fragile] \frametitle{The Update Gate}
\begin{center}
\grudiagram{1}{0}
\end{center}
\begin{itemize}
\item The update gate decides the mix, one value per memory slot: $z_t = \sigma(W_{zx}x_t + W_{zh}h_{t-1} + b_z)$
\item $h_t = z_t \odot h_{t-1} + (1 - z_t) \odot \tilde{h}_t$, where $\odot$ is a slot-by-slot product
\item $z_t$ near 1 copies the old memory; $z_t$ near 0 replaces it with the candidate
\item One slot, $h_{t-1} = 0.8$ and candidate $-0.2$:
  \begin{itemize}
  \item $z_t = 0.9$ gives $0.9 \cdot 0.8 + 0.1 \cdot (-0.2) = 0.70$, mostly kept
  \item $z_t = 0.1$ gives $0.1 \cdot 0.8 + 0.9 \cdot (-0.2) = -0.10$, mostly replaced
  \end{itemize}
\end{itemize}
\end{frame}


%%%%%%%%%%%%%%%%%%%%%%%%%%%%%%%%%%%%%%%%%%%%%%%%%%%
\begin{frame}[fragile] \frametitle{The Reset Gate}
\begin{center}
\grudiagram{1}{1}
\end{center}
\begin{itemize}
\item The reset gate sits on the path from the old state into the candidate: $r_t = \sigma(W_{rx}x_t + W_{rh}h_{t-1} + b_r)$
\item Candidate: $\tilde{h}_t = \tanh(W_{cx}x_t + W_{ch}(r_t \odot h_{t-1}) + b_c)$
\item $r_t$ near 1: the candidate sees the full history, like a plain RNN cell
\item $r_t$ near 0: the candidate is built from the new input alone, a fresh start for the proposal
\item Two different jobs: the update gate decides how much to \emph{replace}; the reset gate decides what the proposal may \emph{look at}
\end{itemize}
\end{frame}


%%%%%%%%%%%%%%%%%%%%%%%%%%%%%%%%%%%%%%%%%%%%%%%%%%%
\begin{frame}[fragile] \frametitle{GRU in Equations}
\begin{align*}
z_t =& \sigma(W_{zx} x_t + W_{zh} h_{t-1} + b_z) \\
r_t =& \sigma(W_{rx} x_t + W_{rh} h_{t-1} + b_r) \\
\tilde{h}_t =& \tanh(W_{cx} x_t + W_{ch} (r_t \odot h_{t-1}) + b_c) \\
h_t =& z_t \odot h_{t-1} + (1 - z_t) \odot \tilde{h}_t
\end{align*}
\begin{itemize}
\item Three blocks of weights, each the size of a plain RNN's: about three times its parameters at the same width
\item One state is carried from step to step: $h_t$
\item As before, the same weights are used at every step
\end{itemize}
\end{frame}


%%%%%%%%%%%%%%%%%%%%%%%%%%%%%%%%%%%%%%%%%%%%%%%%%%%
\begin{frame}[fragile] \frametitle{What a GRU Gives and What Remains}
\textbf{Gives}
\begin{itemize}
\item Can hold a value across many steps: keep $z_t$ near 1
\item Can ignore an input, or restart its context with the reset gate
\item A short direct path for memory, and for the gradient on its way back
\item One state and fewer parameters than an LSTM: quick to train
\end{itemize}
\textbf{Still an issue}
\begin{itemize}
\item Memory and output are the same vector $h_t$: whatever is remembered is also shown
\item Keep and replace are tied: $z_t$ and $1 - z_t$ cannot be set separately
\item Gates ease the vanishing-gradient problem; they do not remove it
\item Still one step at a time, so very long sequences are slow
\end{itemize}
\end{frame}
